% AI-written, AI cross-reviewed; no human sign-off; not a claim that RH is proved
% Atlas unit P033. Local manuscript v3, item 7.10, heading on printed PDF p. 78.
% UpperTheory.V.ZK.sThetaOf
% paper2/inputs_one.tex:434–467. Verbatim excerpt; SHA-256 b6ab11f1e7f123eece8bab4d10978250dec6a775af23cbfb638e49809b38317f.
% Original TeX uses the macros and cross-reference keys of the source manuscript.

\begin{theorem}[The determinant evaluation on scalar prisms]\label{pn:P033}
Assume the chain, prism, prism composition and basic-cycle identities,
the transfer of scalar functionals and the scalar cocycle formula.
For any finite field \(k\), positive integer \(q\), \(e\ge0\),
primitive \(q\)-th root \(\zeta\in k^\times\), and character
\(v:k^\times\to\mathbb Z/q\), there is an additive evaluating
family \(\theta\) of degree \(2e+3\). For every finite \(k\)-vector
space \(V\), every \(x\in\mathbb N\), and \(g\in\mathrm{GL}(V)\),
\[
 \theta_V\!\left(P_g((j\mapsto\zeta^{xj}I_V)_*z)\right)
 = x^{e+1}v(\det g).
\]
Here \(P_g\) is the conjugation prism, and \(z\) is the specified
degree-\(2e+2\) cyclic bar cycle.
\end{theorem}\nodestatus{P033}
\begin{proof}
For \(q=1\) all evaluations are zero. Otherwise \(q\mid |k^\times|\)
and \(|k|=1\) in \(\mathbb Z/q\). Transfer the supplied scalar
cocycle along the projective lines of \(V\). The resulting family
kills boundaries, is additive on direct sums and recovers the given
functional on scalar cycles in dimension one, by the transfer hypothesis.
For fixed \(x,V\), the scalar action commutes with every \(g\).
Prism composition modulo boundaries thus makes the left hand side
a homomorphism from \(\mathrm{GL}(V)\) to an abelian group. It
vanishes on transvections: conjugating \(1+tE_{ij}\) by a diagonal
matrix with \(i\)-th entry \(b\ne0,1\) multiplies \(t\) by \(b\).
Taking differences realizes any transvection since \(b-1\) is
invertible; such a \(b\) exists because \(q>1\). Row elimination,
including signed interchanges expressed as transvections, reduces
the value to that on diagonal matrices. The direct-sum law adds the
one-dimensional evaluations, giving \(x^{e+1}v(\det g)\). The
zero-dimensional case follows from the trivial action. This is the
determinant calculation in Blueprint Z2.K2~\bploc{sec:P-r4-K2}.
\end{proof}
